% !TeX root = ../../Book2.tex
% 纲要标题：### A.4 有限域与有理数上的精确算例
% 纲要位置：计划纲要.md，第 1770 行。

% 纲要要点（供目录核对）：
%
% * 有限域和有理数上的可复核例题；形式多项式与有限点函数、目标余核与来源坐标、环生成数与域维数、分式域可逆与原环满射分别核对；有理相似只须非零行列式，先证明双侧逆再给逆公式；正式证明全部逐点秩仍不能恢复重数
% #### A.4.1 有限域上的多项式矩阵
% #### A.4.2 有理数上的循环基与相似证书
%

\section{有限域与有理数上的精确算例}
\label{b2:sec:A04}

下面的两个算例分别计算多项式矩阵的余核和有理矩阵的相似类。前者使用左右两个独立的换基，后者要求同一线性空间在定义域与值域采用同一组新基。把这两种等价关系混在一起，即使矩阵乘法都算对了，也会回答另一个问题。

\subsection{有限域上的多项式矩阵}
\label{b2:subsec:A04:01}

取 \(R=\mathbb F_5[t]\)，考虑
\[
B=\begin{pmatrix}t^2+1&t\\t&t^2\end{pmatrix}.
\]
所有系数在 \(\mathbb F_5\) 中运算，但 \(t\) 仍是形式多项式变量，不把它代成域中某个数。\(\mathbb F_5[t]\) 与有限集合 \(\mathbb F_5\) 上的函数环不是同一个对象：非零多项式 \(t^5-t\) 在每个 \(a\in\mathbb F_5\) 处的值都为零，所以逐点求值会把不同多项式变成同一个函数。下面的变换均在多项式环中进行。先作 \(C_1\leftarrow C_1-tC_2\)，得到
\[
\begin{pmatrix}1&t\\t-t^3&t^2\end{pmatrix}.
\]
再作 \(R_2\leftarrow R_2+(t^3-t)R_1\)，清除第一列下项，最后作 \(C_2\leftarrow C_2-tC_1\)。这给出
\begin{equation}\label{b2:eq:A-polynomial-certificate}
P=\begin{pmatrix}1&0\\t^3-t&1\end{pmatrix},\qquad
Q=\begin{pmatrix}1&-t\\-t&1+t^2\end{pmatrix},\qquad
D=PBQ=\begin{pmatrix}1&0\\0&t^4\end{pmatrix}.
\end{equation}
这里 \(\det P=\det Q=1\)，所以换基始终在 \(\mathbb F_5[t]\) 上可逆。独立检查中，一阶子式生成单位理想，因为 \((t^2+1)-t\cdot t=1\)；行列式为 \(t^4\)。由\cref{b2:thm:smith-uniqueness}，不变因子确为 \(1,t^4\)。

\ProseParagraphBreak
由于 \(Q\) 可逆，\(PB(R^2)=D(R^2)\)，所以 \(P\) 诱导同构
\[
\operatorname{coker}B\xrightarrow{\sim}\operatorname{coker}D,
\qquad x+B(R^2)\longmapsto Px+D(R^2).
\]
对 \(D=\operatorname{diag}(1,t^4)\)，取第二坐标模去 \(t^4\) 又给出 \(\operatorname{coker}D\cong R/(t^4)\)，与\cref{b2:prop:smith-cokernel}一致。两个同构的复合在原自由模上由映射
\[
R^2\longrightarrow R/(t^4),\qquad
\binom{u}{v}\longmapsto (t^3-t)u+v+(t^4)
\]
给出；其核恰是 \(B(R^2)\)。标准余核中的生成元 \(e_2+D(R^2)\) 经逆同构回到
\[
P^{-1}e_2+B(R^2).
\]
本例恰有 \(P^{-1}e_2=e_2\)，所以原坐标中 \(e_2\) 的陪集就生成这个循环模；一般恢复生成元时须实际计算 \(P^{-1}e_i\)，不能略掉这一换基。\(Q\) 则改变来源自由模中的关系坐标，其可逆性使关系列表的重组不改变 \(B(R^2)\)。商模作为 \(\mathbb F_5\) 向量空间的一组基为
\[
\overline e_2,\quad t\overline e_2,\quad
t^2\overline e_2,\quad t^3\overline e_2.
\]
乘 \(t\) 的算子在这组基下是 \(C(t^4)\)，其最小多项式为 \(t^4\)。作为 \(R\) 模，只需一个生成元 \(\overline e_2\)，因为 \(R\) 系数已允许乘上 \(t\) 的任意多项式；作为 \(\mathbb F_5\) 向量空间，却需要上述四个基元素。这两个生成数使用的系数环不同。

\ProseParagraphBreak
若只计算 \(B\) 在分式域 \(\mathbb F_5(t)\) 上的秩，会得到二：行列式 \(\det B=t^4\) 在该域中是单位，所以矩阵在那里可逆。但 \(t^4\) 不是 \(R=\mathbb F_5[t]\) 中的单位；非零常数才是这个多项式环的单位。因此不能据分式域上的可逆性推出原映射在 \(R\) 上满射，原余核仍是非零的 \(R/(t^4)\)。

\subsection{有理数上的循环基与相似证书}
\label{b2:subsec:A04:02}

在 \(\mathbb Q^3\) 上取算子矩阵
\[
T=\begin{pmatrix}
\tfrac12&-\tfrac14&3\\
1&-\tfrac12&2\\
0&1&0
\end{pmatrix},\qquad v=\begin{pmatrix}1\\0\\0\end{pmatrix}.
\]
逐次乘法给出
\[
Tv=\begin{pmatrix}\tfrac12\\1\\0\end{pmatrix},\qquad
T^2v=\begin{pmatrix}0\\0\\1\end{pmatrix},\qquad
T^3v=\begin{pmatrix}3\\2\\0\end{pmatrix}=2v+2Tv.
\]
前三个向量组成列矩阵
\[
S=(v\ Tv\ T^2v)=
\begin{pmatrix}1&\tfrac12&0\\0&1&0\\0&0&1\end{pmatrix},\qquad \det S=1.
\]
因此它们是一组循环基。在 \(\mathbb Q\) 上，合法换基只要求 \(\det S\ne0\)，本例行列式等于一不是相似变换所必需的条件；非零有理数，如二或 \(1/3\)，也能作为换基矩阵的行列式。这与整数 Smith 变换要求行列式为 \(\pm1\) 不同。记 \(f(t)=t^3-2t-2\)，则其友矩阵为
\[
C(f)=\begin{pmatrix}0&0&2\\1&0&2\\0&1&0\end{pmatrix}.
\]
三列所表示的等式合并为如下相似证书：
\begin{equation}\label{b2:eq:A-rational-similarity}
TS=SC(f),\qquad \det S=1,
\qquad S^{-1}TS=C(f).
\end{equation}
由相似证书，特征多项式为 \(\chi_T(t)=f(t)\)。最小多项式也能直接核验：次数至多为二的非零多项式不可能零化 \(v\)，否则 \(v,Tv,T^2v\) 线性相关；而 \(f(T)v=0\)，与 \(T\) 交换后得到 \(f(T)T^jv=0\)，所以 \(f(T)\) 零化整个空间。最小多项式因而次数为三，并整除首一三次多项式 \(f\)，故 \(\mu_T(t)=\chi_T(t)=f(t)\)，有理标准形只有这一个块。

\ProseParagraphBreak
还可以用同一证书求逆。先将 \(T^3-2T-2I=0\) 改写为
\[
T(T^2-2I)=2I.
\]
\(T^2-2I\) 是 \(T\) 的多项式，与 \(T\) 交换，所以另一顺序的乘积也为 \(2I\)。因此 \(\dfrac12(T^2-2I)\) 已被证明是 \(T\) 的双侧逆，才得到
\begin{equation}\label{b2:eq:A-rational-inverse}
T^{-1}=\tfrac12(T^2-2I).
\end{equation}
全程只使用有理数运算，没有求近似特征值，也不必先把 \(f\) 的根写出来。

\ProseParagraphBreak
\eqref{b2:eq:A-polynomial-certificate}中的 \(P,Q\) 是两个自由模的独立换基，保存矩阵给出的余核；\eqref{b2:eq:A-rational-similarity}中的 \(S^{-1}TS\) 是同一空间的换基，保存算子的相似类。若把后一种换基写成左右乘法，左右矩阵必须互逆，即 \(P=S^{-1},Q=S\)。独立的行列操作一般不保持特征多项式。判断余核同构或线性方程的可解性时使用前一种证书，判断算子相似性时使用后一种证书。

\ProseParagraphBreak
逐点求值能检查某些性质，却不能恢复多项式矩阵的全部不变因子。即使在所有扩域中知道每个点的秩，零点处的重数仍可能没有被记录。

\begin{proposition}\label{b2:prop:A-pointwise-ranks-lose-multiplicity}
设 \(k\) 为任意域，\(r,s\) 为不同的正整数，\(R=k[t]\)，并令
\[
B_r=\operatorname{diag}(1,t^r),\qquad
B_s=\operatorname{diag}(1,t^s)
\]
表示从 \(R^2\) 到 \(R^2\) 的 \(R\) 线性映射。对每个扩域 \(K/k\) 和每个 \(a\in K\)，特殊化矩阵 \(B_r(a)\)、\(B_s(a)\) 的秩相同；但
\[
\operatorname{coker}B_r\cong k[t]/(t^r),\qquad
\operatorname{coker}B_s\cong k[t]/(t^s)
\]
的 \(k\) 维数分别为 \(r,s\)，所以这两个余核作为 \(R\) 模不同构。特别地，当 \(k=\mathbb F_5\)、\(r=4,s=8\) 时，两矩阵在每个 \(a\in\mathbb F_5\) 处的取值甚至完全相同。
\end{proposition}
\begin{proof}
对任意 \(K/k\) 和 \(a\in K\)，若 \(a=0\)，两个矩阵都为 \(\operatorname{diag}(1,0)\)，秩为一；若 \(a\ne0\)，正整数次幂 \(a^r,a^s\) 均非零，两个矩阵的秩都为二。因此全部逐点秩相同。

在 \(R\) 上，第一坐标的像为整个 \(R\)，第二坐标的像分别为 \(t^rR,t^sR\)。取第二坐标的商，便得到所列余核同构。由首一多项式的带余除法，\(R/(t^r)\) 的 \(k\) 基为 \(1,t,\ldots,t^{r-1}\) 的商类，而 \(R/(t^s)\) 有 \(s\) 个相应基元素。任何 \(R\) 模同构也必须保持 \(k\) 标量作用，因而给出 \(k\) 线性同构；维数不同排除这种可能。

最后在 \(\mathbb F_5\) 中，\(a=0\) 时 \(a^4=a^8=0\)；乘法群 \(\mathbb F_5^\times\) 的阶为四，所以 \(a\ne0\) 时 \(a^4=1\)，进而 \(a^8=1\)。故 \(B_4(a)=B_8(a)\) 对每个 \(a\in\mathbb F_5\) 成立，而前段已证明这两个余核不同构。这同时区分了有限点上的矩阵函数和多项式矩阵本身。
\end{proof}
